\section{Echo Chamber Analyse}



Echo chamber effekt
effekt erstmal anhand von literatur erklären mathematisch zeigen, wie sich der effekt auswirkt einmal bei der utitly theory fomel und einmal bei der sigmoid formel und unterschied kurz erklären was evtl. besser geeigent wäre 


%%%%%%%%%%%%%%%
Echo-Chamber-Effekt im Kontext kollaborativer IDS-Systeme definieren

Relevante Literatur (z.B. Debar und Wespi , Heigl et al. )

Übertragung auf das vorliegende P2P-Szenario: Mehrheit fehlerhafter Knoten beeinflusst Aggregat
%%%%%%%%%%%%%%%
Szenario formalisieren: M Knoten melden fälschlicherweise Attack, K Knoten korrekt Benign

Zeige anhand der Aggregationsformel aus Kapitel 6.3, wie das Gewichtsverhältnis M/K den Aggregat-Wert verzerrt
%%%%%%%%%%%%%%%%
Berechnung des aggregierten Outputs für das obige Szenario mit Formel (1)

Durch den frühen starken Anstieg: bereits wenige fehlerhafte Knoten mit mittlerer Detektionszahl kippen den Aggregat
%%%%%%%%%%%%
Effekt unter Sigmoid-Funktion:

Selbes Szenario mit Formel (2)

Durch Wendepunkt n0: fehlerhafte Knoten mit wenigen Detections haben geringeren Einfluss, starke Ausreißer jedoch ähnlich wie Utility Theory

%%%%%%%%%%%%
Tabellarischer Vergleich: Anfälligkeit beider Funktionen gegenüber Echo-Chamber unter verschiedenen M/K-Verhältnissen

Aussage: Keine der Funktionen löst das Problem vollständig -> Verweis auf Kapitel 8.3 als offene Forschungsfrage

Mögliche Gegenmaßnahme als Ausblick: reputationsbasierte Gewichtung, Byzantine-fault-tolerante Aggregation